\subproblem{Интерференция звука}{3.0}

Две бесконечные жесткие плоскости (стены) пересекаются под прямым углом. Из бесконечности падает плоская звуковая волна, волновой вектор которой направлен точно по биссектрисе угла между плоскостями. Амплитуда скорости молекул воздуха в падающей волне равна $v$, а длина волны составляет $\lambda$. Стены идеально отражающие (отражения происходят без потерь энергии); дифракционными эффектами на расстоянии порядка $\lambda$ пренебречь.

\begin{task}[type=subproblem]
    \item Найдите максимально возможную скорость движения воздуха в данной системе.
    \item В каких точках пространства достигается эта максимальная скорость?
\end{task}
